\chapter*{پیش‌گفتار}
\addcontentsline{toc}{chapter}{پیش‌گفتار}

این جزوه برای درس «نظریه زبان‌ها و ماشین‌ها» نوشته شده است. پایه آن اسلایدهای درس دکتر مهدی دولتی در دانشگاه صنعتی شریف است. ترتیب مطالب تقریباً همان ترتیب جلسه‌های درس است، اما مطالب هر جلسه به شکل یک فصل منظم با تعریف، قضیه، اثبات و مثال بازنویسی شده‌اند تا جزوه مثل یک کتاب خوانده شود، نه مثل مجموعه‌ای از اسلایدهای پشت سر هم.

\section*{این جزوه برای چه کسی است؟}
برای خواندن این جزوه لازم نیست از قبل چیزی درباره اتوماتا بدانید. فقط کمی آشنایی با منطق و مجموعه‌ها کافی است و همان‌ها هم در فصل‌های ۱ و ۲ مرور شده‌اند. هر مفهوم جدید اول با یک توضیح شهودی و مثال معرفی می‌شود و بعد تعریف دقیق آن می‌آید.

\section*{درباره اثبات‌ها}
هر جا در اسلایدها فقط «ایده اثبات» آمده بود یا جای اثبات خالی مانده بود (مثلاً با جمله‌هایی مثل «\lr{Formal correctness proof?}» یا «\lr{Exercise}»)، اثبات کامل آن در این جزوه نوشته شده است. اثبات‌ها با زبانی روان نوشته شده‌اند: اول می‌گوییم چه چیزی را می‌خواهیم نشان دهیم، بعد قدم‌به‌قدم جلو می‌رویم. اگر اثباتی طولانی است، قبل از آن یک «ایده اثبات» آمده تا نقشه کلی را ببینید.

\section*{راهنمای کادرها}
مطالب جزوه در کادرهای رنگی جدا شده‌اند تا راحت‌تر مرور شوند:

\begin{definition}[نمونه]
کادر آبی برای \textbf{تعریف‌ها} است. هر چیزی که باید دقیق بدانید در این کادرهاست.
\end{definition}
\begin{theorem}[نمونه]
کادر بنفش برای \textbf{قضیه‌ها} است. کادرهای \textbf{لم} و \textbf{نتیجه} هم رنگ‌های نزدیک به این دارند.
\end{theorem}
\begin{proofidea}
این کادر خط‌چین، نقشه کلی اثبات را قبل از جزئیات نشان می‌دهد.
\end{proofidea}
\begin{proof}
اثبات‌ها با یک نوار خاکستری در کنار متن مشخص شده‌اند و با علامت $\blacksquare$ تمام می‌شوند.
\end{proof}
\begin{example}[نمونه]
مثال‌ها با نوار سبز مشخص شده‌اند.
\end{example}
\begin{note}
نکته‌های مهم و دام‌های رایج در این کادر نارنجی آمده‌اند.
\end{note}

تعریف‌ها، قضیه‌ها و مثال‌های هر فصل با یک شمارنده مشترک شماره‌گذاری شده‌اند؛ مثلاً «قضیه ۴.۳» یعنی سومین کادر شماره‌دار فصل ۴.

\section*{نمادها و قراردادها}
\begin{itemize}
\item اصطلاح‌های مهم در اولین بار با معادل انگلیسی آمده‌اند، مثل \term{زبان منظم}{Regular Language}. در انتهای جزوه یک واژه‌نامه کوتاه هم هست.
\item $\eps$ رشته تهی است، $\Sigma$ الفبای ورودی است و $\blank$ نماد خانه خالی نوار ماشین تورینگ است.
\item $\enc{O}$ کدگذاری (نمایش به صورت رشته) شیء $O$ است.
\item در شکل‌ها، حالت شروع با یک پیکان ورودی و حالت پذیرش با دایره دوخطی نشان داده شده است.
\end{itemize}

\section*{منابع}
مرجع اصلی درس کتاب \lr{Introduction to the Theory of Computation} نوشته مایکل سیپسر است. بخش‌هایی از درس (فرم‌های نرمال، گرامرهای منظم، \lr{CYK} و \lr{PDA}های قطعی) از کتاب \lr{An Introduction to Formal Languages and Automata} نوشته پیتر لینز است. مثال‌های فصل لم تزریق هم از اسلایدهای \lr{Costas Busch} آمده‌اند.

\section*{ساختار جزوه}
جزوه پنج بخش دارد که تقریباً به ترتیب از ماشین ضعیف‌تر به قوی‌تر جلو می‌روند:
\begin{enumerate}
\item \textbf{پیش‌نیازها:} منطق، مجموعه، تابع، گراف، رشته، زبان و روش‌های اثبات.
\item \textbf{زبان‌های منظم:} ماشین‌های متناهی قطعی و غیرقطعی، عبارات منظم، گرامرهای منظم و لم تزریق.
\item \textbf{زبان‌های مستقل از متن:} گرامرهای مستقل از متن، فرم‌های نرمال، الگوریتم \lr{CYK}، ماشین‌های پشته‌ای و لم تزریق برای این زبان‌ها.
\item \textbf{محاسبه‌پذیری:} ماشین تورینگ، تز چرچ–تورینگ، تصمیم‌ناپذیری، کاهش و مسئله تناظر پست.
\item \textbf{پیچیدگی:} کلاس‌های $\cP$ و $\cNP$ و $\cNP$-تمامیت.
\end{enumerate}
